\documentclass[11pt]{article}

% ---------- 常用宏包 ----------
\usepackage[a4paper,margin=1in]{geometry}
\usepackage{amsmath,amssymb,amsthm}
\usepackage{graphicx}
\usepackage{booktabs}
\usepackage{hyperref}
\usepackage{xcolor}

\hypersetup{
  colorlinks=true,
  linkcolor=blue!50!black,
  urlcolor=blue!50!black,
  citecolor=blue!50!black
}

\title{A Brand New \LaTeX{} Document}
\author{Your Name}
\date{\today}

\begin{document}

\maketitle

\begin{abstract}
  在这里写摘要。摘要通常用一两段话概括研究问题、方法和主要结论。
\end{abstract}

\section{引言}
在这部分介绍研究背景与动机。行内公式示例：$E = mc^2$，行间公式示例如下：
\begin{equation}
  \int_{-\infty}^{\infty} e^{-x^{2}}\,\mathrm{d}x = \sqrt{\pi}.
  \label{eq:gauss}
\end{equation}

\section{方法}
引用公式~\eqref{eq:gauss} 时可以使用 \verb|\eqref| 命令。

\begin{itemize}
  \item 第一点结论；
  \item 第二点结论。
\end{itemize}

\section{结论}
总结全文的主要发现。

\end{document}
